\slajdid{10-s017}
\begin{frame}[label=10-s017]
\gusttekst\setlength{\abovedisplayskip}{3pt}\setlength{\belowdisplayskip}{3pt}\setlength{\abovedisplayshortskip}{2pt}\setlength{\belowdisplayshortskip}{2pt}Ostaje da odredimo pri kojem ulaznom naponu $V_{I1}$ tranzistor ulazi u zasićenje
\[V_O=V_{CC}-\beta_F\frac{R_C}{R_B}(V_{I1}-V_{BE})=V_{CES}\]
Tada je
\[V_{I1}=\frac{R_B}{\beta_F R_C}(V_{CC}-V_{CES})+V_{BE}\]
\begin{columns}[c,onlytextwidth]\column{.34\textwidth}\centering\begin{tikzpicture}[x=.95cm,y=.8cm,font=\sitneoznake,>=Latex]\draw[->](0,0)--(3.9,0)node[below]{$V_I$};\draw[->](.2,-.2)--(.2,3.1)node[left]{$V_O$};\node[left]at(.2,2.7){$V_{CC}$};\draw(.2,2.7)--(.85,2.7)--(2.3,.18)--(3.25,.18);\draw[dashed](.85,2.7)--(.85,0)node[below]{$V_{\gamma T}$};\draw[dashed](.2,.18)--(2.3,.18);\node[left]at(.2,.18){$V_{CES}$};\draw(3.25,.18)--(3.25,0)node[below]{$V_{CC}$};\node[right]at(1.55,1.6){a};\draw(2.3,.18)--(2.3,0)node[below]{$V_{I1}$};\end{tikzpicture}
\column{.08\textwidth}\tikz\draw[->,DEplava](0,0)--(.8,0);\column{.56\textwidth}
\[V_{OH}=V_{CC}\qquad V_{OL}=V_{CES}\]
\[V_{IL}=V_{\gamma T}\qquad V_{IH}=V_{I1}\]
\[NM_{LMS}=V_{IL}-V_{OL}=V_{\gamma T}-V_{CES}\]
\[NM_{HMS}=V_{OH}-V_{IH}=V_{CC}-V_{I1}\]
\end{columns}
\end{frame}
